% ==============================================================================
% DreamZero: World Action Models are Zero-shot Policies
% Paper Review Presentation — WoRV / MaumAI
% ==============================================================================

% ==============================================================================
% Paper2PR — Shared Beamer Preamble
% ==============================================================================
% Paper-agnostic preamble for all paper review presentations.
% Colors and environments match Quarto clean-academic.scss theme.
% ==============================================================================

\documentclass[aspectratio=169, 10pt]{beamer}

% --- Theme ---
\usetheme{default}
\usecolortheme{default}
\setbeamertemplate{navigation symbols}{}
\setbeamertemplate{footline}[frame number]
\setbeamertemplate{itemize items}[circle]
\setbeamertemplate{enumerate items}[default]

% --- Colors (matching Quarto clean-academic.scss) ---
\usepackage{xcolor}
\definecolor{primaryblue}{HTML}{012169}
\definecolor{primarygold}{HTML}{B9975B}
\definecolor{highlightyellow}{HTML}{F2A900}
\definecolor{lightbg}{HTML}{E8EDF5}
\definecolor{jet}{HTML}{1A1A1A}
\definecolor{positivegreen}{HTML}{15803D}
\definecolor{negativered}{HTML}{B91C1C}

\setbeamercolor{title}{fg=primaryblue}
\setbeamercolor{subtitle}{fg=primarygold}
\setbeamercolor{frametitle}{fg=primaryblue}
\setbeamercolor{structure}{fg=primaryblue}
\setbeamercolor{normal text}{fg=jet}
\setbeamercolor{itemize item}{fg=primaryblue}
\setbeamercolor{itemize subitem}{fg=primarygold}
\setbeamercolor{enumerate item}{fg=highlightyellow}

% --- Fonts (XeLaTeX) ---
\usepackage{fontspec}
% Source Sans Pro (static OTF from TexLive — explicit path for reliability)
\setsansfont{SourceSansPro}[
  Path = /usr/local/texlive/2024/texmf-dist/fonts/opentype/adobe/sourcesanspro/,
  UprightFont = SourceSansPro-Regular.otf,
  BoldFont = SourceSansPro-Bold.otf,
  ItalicFont = SourceSansPro-RegularIt.otf,
  BoldItalicFont = SourceSansPro-BoldIt.otf,
  Scale = 1.0
]
\usefonttheme{professionalfonts}

% --- Math ---
\usepackage{amsmath}
\usepackage{amssymb}
\usepackage{bm}
\usepackage{mathtools}

% --- Graphics ---
\usepackage{graphicx}
\usepackage{tikz}
\usepackage{booktabs}
\usepackage{multicol}
\usepackage{hyperref}

% --- tcolorbox for custom environments ---
\usepackage[most]{tcolorbox}

% methodbox — Blue-bordered for technical details
\newtcolorbox{methodbox}[1][]{%
  colback=primaryblue!6,
  colframe=primaryblue,
  coltitle=white,
  fonttitle=\bfseries,
  left=8pt, right=8pt, top=6pt, bottom=6pt,
  boxrule=0pt, leftrule=4pt,
  arc=4pt,
  title={#1}
}

% keybox — Gold-bordered for key insights
\newtcolorbox{keybox}[1][]{%
  colback=primarygold!12,
  colframe=primarygold,
  coltitle=white,
  fonttitle=\bfseries,
  left=8pt, right=8pt, top=6pt, bottom=6pt,
  boxrule=0pt, leftrule=4pt,
  arc=4pt,
  title={#1}
}

% highlightbox — Yellow-bordered for notable findings
\newtcolorbox{highlightbox}[1][]{%
  colback=highlightyellow!8,
  colframe=highlightyellow,
  coltitle=white,
  fonttitle=\bfseries,
  left=8pt, right=8pt, top=6pt, bottom=6pt,
  boxrule=0pt, leftrule=4pt,
  arc=4pt,
  title={#1}
}

% assumptionbox — Gold full-border for assumptions
\newtcolorbox{assumptionbox}[1][]{%
  colback=primarygold!8,
  colframe=primarygold,
  coltitle=white,
  fonttitle=\bfseries,
  left=8pt, right=8pt, top=6pt, bottom=6pt,
  boxrule=2pt,
  arc=4pt,
  title={#1}
}

% resultbox — Gold-bordered with shadow for main results
\newtcolorbox{resultbox}[1][]{%
  colback=primarygold!15,
  colframe=primarygold,
  coltitle=white,
  fonttitle=\bfseries,
  left=8pt, right=8pt, top=6pt, bottom=6pt,
  boxrule=2pt,
  arc=4pt,
  shadow={1pt}{-1pt}{0pt}{black!12},
  title={#1}
}

% eqbox — Subtle blue background for equations
\newtcolorbox{eqbox}{%
  colback=primaryblue!4,
  colframe=primaryblue!4,
  left=8pt, right=8pt, top=4pt, bottom=4pt,
  boxrule=0pt,
  arc=0pt
}

% softbox — Subtle gold italic for remarks
\newtcolorbox{softbox}{%
  colback=primarygold!6,
  colframe=primarygold!6,
  left=8pt, right=8pt, top=4pt, bottom=4pt,
  boxrule=0pt,
  arc=0pt,
  fontupper=\itshape
}

% quotebox — Italic with decorative quote mark
\newtcolorbox{quotebox}{%
  colback=primarygold!8,
  colframe=primarygold!80!black,
  left=16pt, right=8pt, top=6pt, bottom=6pt,
  boxrule=0pt, leftrule=4pt,
  arc=0pt, outer arc=4pt,
  fontupper=\itshape,
  before upper={\textcolor{primarygold!40}{\Large``}\,}
}

% --- Text color commands ---
\newcommand{\hi}[1]{{\bfseries\color{primaryblue}#1}}
\newcommand{\higold}[1]{{\bfseries\color{primarygold}#1}}
\newcommand{\hiyellow}[1]{{\bfseries\color{highlightyellow}#1}}
\newcommand{\higreen}[1]{{\bfseries\color{positivegreen}#1}}
\newcommand{\hired}[1]{{\bfseries\color{negativered}#1}}
\newcommand{\positive}[1]{{\bfseries\color{positivegreen}#1}}
\newcommand{\negative}[1]{{\bfseries\color{negativered}#1}}
\newcommand{\neutral}[1]{{\color{gray}#1}}

% --- Bibliography ---
\usepackage[
  backend=bibtex,
  style=authoryear,
  natbib=true,
  maxcitenames=2
]{biblatex}

% --- Title page template (compact, no internal \vfill) ---
\setbeamertemplate{title page}{
  \begin{center}
    \vspace{1em}
    {\Large\bfseries\color{primaryblue}\inserttitle}

    \vspace{0.4em}
    {\normalsize\color{primarygold}\insertsubtitle}

    \vspace{0.8em}
    {\normalsize\insertauthor}

    \vspace{0.2em}
    {\small\textcolor{primarygold}{\insertinstitute}}
  \end{center}
}

% --- Frame title formatting ---
\setbeamerfont{frametitle}{size=\large, series=\bfseries}
\setbeamertemplate{frametitle}{%
  \vspace{0.3em}
  \insertframetitle
  \vspace{0.1em}
  {\color{primarygold!30}\hrule height 2pt}
}


\addbibresource{../Bibliography_base.bib}

\title{DreamZero: World Action Models\\are Zero-shot Policies}
\subtitle{Ye, Ge, Zheng, Gao, Yu et al.\ (NVIDIA, 2026)}
\author{Yunsung Lee}
\institute{WoRV / MaumAI}
\date{}

\begin{document}

% ==============================================================================
% PART 1: MOTIVATION & BACKGROUND
% ==============================================================================

% --- Slide 1: Title ---
\begin{frame}[plain]
\maketitle
\vfill
\begin{center}
\includegraphics[width=0.85\textwidth,
  height=0.55\textheight,keepaspectratio]{%
  ../Figures/DreamZero/dreamzero-header-v2.pdf}
\end{center}
\end{frame}

% --- Slide 2: The Second Pre-training Paradigm ---
\begin{frame}{The Second Pre-training Paradigm}
\hi{1st Paradigm:} Next-word prediction
$\rightarrow$ Language intelligence (LLMs)

\vspace{0.5em}
\hi{2nd Paradigm:} \higold{Next physical state prediction}
$\rightarrow$ Embodied intelligence (World Models)

\vspace{0.5em}
\begin{quotebox}
2026 will go down as the first year that Large World Models
lay real foundations for robotics.
\end{quotebox}
\hfill --- Jim Fan (NVIDIA)
\end{frame}

% --- Slide 2b: Why Not Just Use VLMs? ---
\begin{frame}{Why Not Just Use VLMs for Robots?}
\begin{itemize}
    \item VLMs are \negative{language-first} ---
          vision is a ``second-class citizen''
    \item Most VLM parameters encode \emph{knowledge},
          not \emph{physics}
    \item Static image--text pretraining
          $\neq$ spatiotemporal understanding
\end{itemize}

\vspace{0.5em}
\begin{softbox}
\higold{Commentary:} ``Language is a bottleneck, a scaffold,
not a foundation'' for physical intelligence.
\end{softbox}
\end{frame}

% --- Slide 3: The Data Problem (1/2) ---
\begin{frame}{The Data Problem in Robotics (1/2)}
\hi{How much data does each domain have?}
\begin{itemize}
    \item \hi{LLMs:} text from \positive{billions of people}
          across history
    \item \hi{Vision:} images from \positive{billions of cameras}
          over decades
    \item \hi{Robotics:} \negative{thousands of hours}
          of teleoperation $\rightarrow$ \negative{massive gap}
\end{itemize}

\vspace{0.5em}
\begin{highlightbox}[In-the-wild vs.\ On-demand Data]
Text and images are ``in-the-wild'' --- generated continuously.\\
Robot data is ``on-demand'' --- requires expensive teleoperation.
\end{highlightbox}
\end{frame}

% --- Slide 3b: The Data Problem (2/2) ---
\begin{frame}{The Data Problem in Robotics (2/2)}
\begin{quotebox}
You can't teleoperate your way to the long tail.
\end{quotebox}
\hfill --- Joel Jang

\vspace{0.5em}
\emph{``Intelligence is a function of experience.''}

\vspace{0.5em}
$\Rightarrow$ Robotics needs \higold{vastly more experience}
to match human-level physical intelligence.
\end{frame}

% --- Slide 4: VLAs and Their Limitations ---
\begin{frame}{The Rise of VLAs and Their Limitations}
\hi{Vision-Language-Action Models (VLAs):}
\begin{itemize}
    \item RT-2, OpenVLA, GR00T, $\pi_0$
          \citep{rt22023arxiv,kim2024openvla,bjorck2025gr00t,black2024pi0}
    \item Extend VLMs $\rightarrow$ predict motor actions
    \item \positive{Strong semantic understanding}:
          ``pick up the red cup''
\end{itemize}

\vspace{0.5em}
\hi{But:}
\begin{itemize}
    \item Trained on \negative{static images}
          $\rightarrow$ lack spatiotemporal priors
    \item Fail on novel \emph{motions}:
          ``untie the shoelace''
    \item Scaling model size alone
          \negative{does not help}
\end{itemize}
\end{frame}

% --- Slide 4b: VLAs — The Evidence ---
\begin{frame}{VLAs: The Evidence}
\begin{center}
\includegraphics[width=0.75\textwidth,
  height=0.45\textheight,keepaspectratio]{%
  ../Figures/DreamZero/eval_bar_chart.pdf}
\end{center}

{\small VLA from-scratch at $\sim$0\% on diverse data}

\vspace{0.5em}
\begin{softbox}
\higold{Commentary:} The fundamental gap --- VLMs encode
``\hi{what} to do'' but not ``\hi{how} to move''.
\end{softbox}
\end{frame}

% --- Slide 5: Video as Dense Representation ---
\begin{frame}{Key Insight: Video as Dense Physical Representation}
\begin{keybox}[Core Idea]
Video diffusion models trained on web-scale data encode:\\
\hi{spatial geometry} $\cdot$ \hi{object dynamics}
$\cdot$ \hi{contact physics} $\cdot$ \hi{motion patterns}
\end{keybox}

\vspace{0.5em}
\begin{itemize}
    \item Every consecutive frame pair
          $\rightarrow$ \positive{dense supervision}
          (vs.\ sparse task-level labels in VLAs)
    \item ``Generating pixels is the only way for humans
          to \emph{verify} that the model truly understands''
          --- Joel Jang
\end{itemize}
\end{frame}

% --- Slide 5b: The Paradigm Shift ---
\begin{frame}{The Paradigm Shift}
\begin{quotebox}
Apes may not have good language models, but they surely
have a robust mental picture of counterfactuals.
\end{quotebox}
\hfill --- Jim Fan

\vspace{0.5em}
\begin{softbox}
\higold{Commentary:} This is the core paradigm shift ---
from discrete task policies to
\hi{continuous world modeling}.
\end{softbox}
\end{frame}

% --- Slide 6: WAM Big Picture ---
\begin{frame}{World Action Models: The Big Picture}
\hi{WAM} = Video Prediction $\times$ Inverse Dynamics Model

\vspace{0.5em}
\begin{eqbox}
\begin{equation*}
\underbrace{P(\mathbf{O}, \mathbf{A} \mid \text{ctx})}_{\text{DreamZero}}
  = \underbrace{P(\mathbf{O} \mid \text{ctx})}_{\text{video pred.}}
  \times \underbrace{P(\mathbf{A} \mid \mathbf{O},
  \text{ctx})}_{\text{IDM}}
\end{equation*}
\end{eqbox}

\vspace{0.3em}
{\small\hi{IDM (Inverse Dynamics Model):} observes a sequence of states
and infers what actions caused those transitions ---
watching video and reverse-engineering motor commands.}

\vspace{0.5em}
\begin{itemize}
    \item Instead of two separate models
          $\rightarrow$ \hi{single end-to-end model}
    \item \hi{DreamZero:} 14B autoregressive WAM built on
          \higold{Wan2.1} video backbone \citep{wan2025wan}
\end{itemize}
\end{frame}

% --- Slide 6b: Headline Results ---
\begin{frame}{DreamZero: Headline Results}
\begin{resultbox}[Key Achievements]
\begin{itemize}
    \item \positive{2$\times$ improvement}
          over state-of-the-art VLAs
    \item \positive{38$\times$ inference speedup}
          $\rightarrow$ real-time at 7Hz
    \item \positive{Cross-embodiment transfer}
          from video-only data
\end{itemize}
\end{resultbox}
\end{frame}

% --- Slide 7: Human Data Scaling ---
\begin{frame}{The Human Data Scaling Pathway}
\begin{quotebox}
Humans are robots that have already been deployed at scale.
\end{quotebox}
\hfill --- Joel Jang

\vspace{0.5em}
\hi{The math:}
\begin{itemize}
    \item 8 billion ``units'' $\times$ 16 waking hrs/day
          $=$ 128B person-hours \emph{per day}
    \item Even capturing 0.1\% $\rightarrow$
          $\approx$ 100M hours of usable video
    \item $\approx$ 150 human lifetimes of experience
\end{itemize}

\vspace{0.5em}
\hi{The vision:}
``Train once on human experience,
deploy everywhere with calibration''
\end{frame}

% --- Slide 7b: Early Signal ---
\begin{frame}{Human Data: Early Signal}
\begin{resultbox}[Cross-Embodiment Transfer]
12 min of human egocentric video\\
$\rightarrow$ \positive{$>$42\% relative improvement}
on unseen robot tasks
\end{resultbox}

\vspace{0.5em}
\begin{softbox}
\higold{Commentary:} This is the sleeper insight ---
human video as the ultimate robot training data.
\end{softbox}
\end{frame}

% ==============================================================================
% PART 2: DREAMZERO ARCHITECTURE
% ==============================================================================

% --- Section Divider ---
\begin{frame}[plain]
\begin{center}
\vfill
{\Huge\hi{How Does DreamZero Work?}}\\[1em]
{\large\higold{Architecture \& Training}}
\vfill
\end{center}
\end{frame}

% --- Slide 8: Architecture Overview ---
\begin{frame}{Architecture Overview}
\begin{center}
\includegraphics[width=0.88\textwidth,
  height=0.50\textheight,keepaspectratio]{%
  ../Figures/DreamZero/dreamzero_model.pdf}
\end{center}

\begin{columns}[T]
\column{0.32\textwidth}
\hi{Inputs:}\\
Visual context (VAE)\\
Language (text encoder)\\
Proprioceptive state
\column{0.32\textwidth}
\hi{Backbone:}\\
Autoregressive DiT\\
{\small(Diffusion Transformer)}\\
with flow matching
\column{0.32\textwidth}
\hi{Outputs:}\\
Joint video frames\\
+ action chunks
\end{columns}
\end{frame}

% --- Slide 9: Why AR? ---
\begin{frame}{Why Autoregressive? AR vs.\ Bidirectional}
\begin{columns}[T]
\column{0.48\textwidth}
\hi{Bidirectional problems:}
\begin{itemize}
    \item Fixed-length sequences
          $\rightarrow$ \negative{video subsampling}
    \item Distorts native FPS
          $\rightarrow$ breaks alignment
    \item Especially harmful mid-task
\end{itemize}

\vspace{0.3em}
\includegraphics[width=\textwidth,
  height=0.28\textheight,keepaspectratio]{%
  ../Figures/DreamZero/bidiraug.pdf}

\column{0.48\textwidth}
\hi{Autoregressive advantages:}
\begin{itemize}
    \item Preserves \positive{native frame rate}
    \item \positive{KV-cache} for 3--4$\times$ speedup
    \item Natural modality alignment
    \item Variable-length training
\end{itemize}
\end{columns}
\end{frame}

% --- Slide 9b: AR — Commentary ---
\begin{frame}{AR vs.\ Bidirectional: Commentary}
\begin{softbox}
\higold{Commentary:} AR wins on both quality AND speed ---
a rare ``have your cake and eat it too'' situation.
\end{softbox}

\vspace{0.5em}
\begin{itemize}
    \item \hi{Speed:} KV-cache gives 3--4$\times$ speedup ---
          bidirectional cannot cache
    \item \hi{Quality:} native frame rate preserves the temporal
          signal that the IDM needs
\end{itemize}
\end{frame}

% --- Slide 10: Chunk-wise Generation ---
\begin{frame}{Chunk-wise Generation}
\begin{methodbox}[Chunking Strategy]
\begin{itemize}
    \item Video: \hi{$K=2$ latent frames} per chunk
    \item Action horizon: \hi{$H=48$} (AgiBot, 30Hz)
          or \hi{$H=24$} (DROID, 15Hz)
    \item Each chunk $\approx$ \higold{1.6 seconds}
    \item Max 4 chunks $\rightarrow$ \higold{6.6 seconds} context
\end{itemize}
\end{methodbox}

\vspace{0.5em}
\hi{Key design choices:}
\begin{itemize}
    \item Variable-length training (like LLM token sequences)
    \item Video at 5 FPS, actions at 30Hz (AgiBot)
    \item 33 raw frames $\rightarrow$ 8 latent frames ($4\times2$)
\end{itemize}
\end{frame}

% --- Slide 11: Flow Matching Intuition ---
\begin{frame}{Flow Matching: Intuition}
\begin{keybox}[What is Flow Matching?]
Learn a \hi{velocity field} that transports noise to data
along straight paths --- simpler and faster than
standard diffusion.
\end{keybox}

\vspace{0.3em}
{\small\emph{Analogy:} Think of a GPS that plots a
straight-line route from ``lost'' (noise) to ``home'' (data).
Flow matching learns that route; standard diffusion
takes a random walk.}

\vspace{0.5em}
\begin{itemize}
    \item At $t=0$: pure noise. At $t=1$: clean data.
    \item The model predicts the \hi{velocity}
          pointing from noise to data.
    \item \hi{Shared timestep} $t_k$ for video and action
          $\rightarrow$ aligned convergence.
\end{itemize}
\end{frame}

% --- Slide 11b: Flow Matching Equations ---
\begin{frame}{Flow Matching: Equations}
\begin{eqbox}
\hi{Noisy interpolation:}
\begin{equation*}
\mathbf{z}_{t_k}^k = t_k \mathbf{z}_1^k
  + (1-t_k)\mathbf{z}_0^k, \quad
\mathbf{a}_{t_k}^k = t_k \mathbf{a}_1^k
  + (1-t_k)\mathbf{a}_0^k
\end{equation*}
where $\mathbf{z}_0^k, \mathbf{a}_0^k
  \sim \mathcal{N}(\mathbf{0}, \mathbf{I})$
and $t_k \in [0,1]$.
\end{eqbox}

\vspace{0.5em}
\begin{eqbox}
\hi{Velocity prediction loss:}
\begin{equation*}
\mathcal{L}(\theta) = \mathbb{E}\Bigg[
  \frac{1}{K}\sum_{k=1}^{K} w(t_k)
  \Big\lVert \mathbf{u}_\theta(\cdots)
  - \mathbf{v}^k \Big\rVert^2\Bigg]
\end{equation*}
\end{eqbox}
\end{frame}

% --- Slide 11c: Flow Matching Notation ---
\begin{frame}{Flow Matching: Notation Guide}
\hi{Key Symbols:}
\begin{itemize}
    \item $\mathbf{u}_\theta$: the DiT model
          (predicts velocity)
    \item $\mathcal{C}_k$: clean context
          from previous chunks
    \item $\mathbf{c}$: language instruction
    \item $\mathbf{q}_k$: proprioceptive state
    \item $w(t_k)$: predefined weight function
\end{itemize}

\vspace{0.5em}
\hi{Teacher forcing:} denoise noisy current chunk
conditioned on \positive{clean} previous chunks
(like LLM training with ground-truth tokens).
\end{frame}

% --- Slide 12: Attention Masking ---
\begin{frame}{Attention Masking Strategy}
\begin{center}
\includegraphics[width=0.78\textwidth,
  height=0.45\textheight,keepaspectratio]{%
  ../Figures/DreamZero/dreamzero_attention_figure.pdf}
\end{center}

\begin{methodbox}[Key Rules]
\begin{itemize}
    \item \hi{Training:} noisy chunk attends ONLY to
          \positive{clean context}
    \item \hi{No cross-attention} between noisy video
          and noisy actions within same chunk
\end{itemize}
\end{methodbox}
\end{frame}

% --- Slide 13: Ground-Truth Injection ---
\begin{frame}{Inference: Ground-Truth Injection}
\begin{keybox}[The Elegant Trick]
\begin{enumerate}
    \item Denoise one chunk at a time (autoregressive)
    \item Execute action chunk on robot
          $\rightarrow$ get \hi{real observation}
    \item \higold{Replace predicted frame with ground-truth}
          in KV cache
    \item This \positive{eliminates error accumulation!}
\end{enumerate}
\end{keybox}

\vspace{0.5em}
\hi{Why this is unique to WAMs:}
\begin{itemize}
    \item Pure video generation: errors compound
    \item WAMs operate in \positive{closed-loop} ---
          reality corrects the model
\end{itemize}
\end{frame}

% --- Slide 13b: GT Injection — Commentary ---
\begin{frame}{Ground-Truth Injection: Commentary}
\begin{softbox}
\higold{Commentary:} WAMs \hi{self-correct via reality}.
No other paradigm gets this for free.
\end{softbox}

\vspace{0.5em}
\begin{itemize}
    \item Closes the sim-to-real gap \emph{by design},
          not by domain randomization
    \item Compare to pure video generation
          (e.g.\ Sora) where errors compound indefinitely
\end{itemize}
\end{frame}

% --- Slide 14: DreamZero-Flash ---
\begin{frame}{DreamZero-Flash: Faster Inference}
\hi{Problem:} 4 denoising steps = 350ms
(too slow for some applications)

\vspace{0.5em}
\hi{Solution:} \higold{Decoupled noise scheduling}
\begin{itemize}
    \item Video: $t_{\text{vid}} = 1 - \eta$,
          $\eta \sim \text{Beta}(7,1)$\\
          $\rightarrow$ biased toward \negative{high-noise}
          (mean $t \approx 0.125$)
    \item Action: $t_{\text{act}} \sim \text{Uniform}[0,1]$
          $\rightarrow$ standard
\end{itemize}

\vspace{0.3em}
Trains actions to denoise from
\emph{predominantly noisy} visual context.
\end{frame}

% --- Slide 14b: Flash Results ---
\begin{frame}{DreamZero-Flash: Results}
\begin{center}
\includegraphics[width=0.60\textwidth,
  height=0.40\textheight,keepaspectratio]{%
  ../Figures/DreamZero/%
  dreamzero-flash_noise_distribution_comparison.pdf}
\end{center}

\begin{resultbox}[Result]
1-step: \positive{74\%} task progress
(only 9\% drop from 4-step at 83\%)
\end{resultbox}
\end{frame}

% --- Slide 15: 38x Speedup ---
\begin{frame}{38$\times$ Inference Speedup Breakdown}
\begin{center}
{\small
\begin{tabular}{lcc}
\toprule
\textbf{Optimization} & \textbf{H100} & \textbf{GB200}\\
\midrule
Baseline & $1\times$ & $1.1\times$ \\
\midrule
\multicolumn{3}{l}{\textit{System-level}} \\
\quad + CFG Parallelism & $1.9\times$ & $1.8\times$ \\
\quad + DiT Caching & $5.5\times$ & $5.4\times$ \\
\midrule
\multicolumn{3}{l}{\textit{Implementation-level}} \\
\quad + torch.compile & $8.9\times$ & $10.9\times$ \\
\quad + Kernel \& Scheduler & $9.6\times$ & $14.8\times$ \\
\quad + Quantization & --- & $16.6\times$ \\
\midrule
\multicolumn{3}{l}{\textit{Model-level}} \\
\quad + DreamZero-Flash & --- & \higold{$38\times$} \\
\bottomrule
\end{tabular}
}
\end{center}

\hi{Final:} 14B at \positive{7Hz} on 2$\times$GB200
(150ms per chunk)
\end{frame}

% ==============================================================================
% PART 3: DATA & TRAINING
% ==============================================================================

% --- Slide 16: Data Philosophy ---
\begin{frame}{Data Philosophy: Diversity Over Repetition}
\begin{resultbox}[Counterintuitive Finding]
Same 500 hours of data:\\
\hi{Diverse} (22 envs, long-tail skills):
  \positive{50\% task progress}\\
\hi{Repetitive} (70 tasks, many reps):
  \negative{33\% task progress}\\
$\rightarrow$ \positive{+51\% relative improvement} from diversity
\end{resultbox}

\vspace{0.5em}
\hi{Why?}
\begin{itemize}
    \item Video prediction inherited from pretraining
    \item Robust IDM requires
          \positive{diverse state-action correspondences}
    \item Repetitive data $\rightarrow$ narrow IDM
\end{itemize}
\end{frame}

% --- Slide 17a: AgiBot Dataset ---
\begin{frame}{AgiBot G1 Dataset}
\begin{columns}[T]
\column{0.50\textwidth}
\hi{Data:}
\begin{itemize}
    \item \higold{500 hours} teleoperation
    \item 7.2K episodes
    \item \higold{22} real-world environments
    \item Avg episode: 4.4 min, $\sim$42 subtasks
\end{itemize}

\column{0.46\textwidth}
\includegraphics[width=\textwidth,
  height=0.45\textheight,keepaspectratio]{%
  ../Figures/DreamZero/agibot_dc_envs.pdf}

{\small 22 diverse environments}
\end{columns}
\end{frame}

% --- Slide 17b: Training Config ---
\begin{frame}{Training Configuration}
\begin{columns}[T]
\column{0.50\textwidth}
\hi{Training:}
\begin{itemize}
    \item Backbone: \hi{Wan2.1-I2V-14B}
    \item 100K steps, batch size 128
    \item \positive{Freeze:} text enc, image enc, VAE
    \item \hi{Update:} DiT + state/action enc/dec
    \item Post-training: 50K steps per task
\end{itemize}

\column{0.46\textwidth}
\includegraphics[width=\textwidth,
  height=0.45\textheight,keepaspectratio]{%
  ../Figures/DreamZero/agibot_stats_3.pdf}

{\small Skill distribution}
\end{columns}
\end{frame}

% --- Slide 18a: Evaluation Setup ---
\begin{frame}{Evaluation Protocol: Setup}
\begin{center}
\includegraphics[width=0.80\textwidth,
  height=0.55\textheight,keepaspectratio]{%
  ../Figures/DreamZero/main_eval_setting.pdf}
\end{center}

Default: \hi{unseen environments}, \hi{unseen objects}\\
Training and evaluation in
\negative{different geographic locations}
\end{frame}

% --- Slide 18b: Evaluation Details ---
\begin{frame}{Evaluation Protocol: Details}
\begin{methodbox}[True Out-of-Distribution]
\begin{itemize}
    \item Seen tasks: 10 tasks $\times$ 8 rollouts
          $\times$ 4 robots = 80 rollouts
    \item Unseen tasks: 10 tasks absent from training
    \item Baselines: GR00T N1, $\pi_{0.5}$
          (from-scratch \& from-pretrained)
\end{itemize}
\end{methodbox}
\end{frame}

% ==============================================================================
% PART 4: RESULTS
% ==============================================================================

% --- Section Divider ---
\begin{frame}[plain]
\begin{center}
\vfill
{\Huge\hi{Does It Actually Work?}}\\[1em]
{\large\higold{Experimental Results}}
\vfill
\end{center}
\end{frame}

% --- Slide 19: Seen Tasks ---
\begin{frame}{Seen Tasks: 2$\times$ Improvement}
\begin{center}
\includegraphics[width=0.78\textwidth,
  height=0.50\textheight,keepaspectratio]{%
  ../Figures/DreamZero/eval_bar_chart.pdf}
\end{center}

\begin{resultbox}
DreamZero: \positive{62.2\%} avg.\ task progress\\
Best VLA (pretrained): 27.4\%
$\rightarrow$ \positive{2.3$\times$}\\
VLA from-scratch: \negative{$\sim$0\%}
\end{resultbox}
\end{frame}

% --- Slide 20: Unseen Tasks ---
\begin{frame}{Unseen Tasks: Zero-shot Generalization}
\begin{center}
\includegraphics[width=0.78\textwidth,
  height=0.50\textheight,keepaspectratio]{%
  ../Figures/DreamZero/unseen_success_rates.pdf}
\end{center}

\begin{resultbox}
DreamZero: \positive{39.5\%} on novel tasks\\
Best VLA (pretrained): 16.3\%
$\rightarrow$ \positive{2.4$\times$}
\end{resultbox}
\end{frame}

% --- Slide 20b: Unseen Tasks Details ---
\begin{frame}{Unseen Tasks: Standout Results}
\hi{Standout tasks:}
\begin{itemize}
    \item Remove hat: \positive{85.7\%}
    \item Shake hands: \positive{59.2\%}
\end{itemize}

\vspace{0.3em}
\hi{VLAs overfit} to pick-and-place
regardless of instruction.

\vspace{0.5em}
\begin{softbox}
\higold{Commentary:} The key result --- DreamZero does
things it was \hi{NEVER trained to do}.
\end{softbox}
\end{frame}

% --- Slide 21: Joint Video-Action ---
\begin{frame}{Joint Video-Action Prediction in Action}
\begin{center}
\includegraphics[width=0.85\textwidth,
  height=0.50\textheight,keepaspectratio]{%
  ../Figures/DreamZero/dreamzero_gen.pdf}
\end{center}

\begin{itemize}
    \item Tight alignment between predicted
          and actual trajectories
    \item Most failures: video generation errors,
          not action extraction
\end{itemize}
\end{frame}

% --- Slide 22: Post-training ---
\begin{frame}{Post-training: Retained Generalization}
\begin{center}
\includegraphics[width=0.78\textwidth,
  height=0.50\textheight,keepaspectratio]{%
  ../Figures/DreamZero/posttrain.pdf}
\end{center}

\begin{itemize}
    \item \positive{Matches or outperforms} pretrained VLAs
    \item Evaluated in \hi{unseen environments}
          even after task-specific training
\end{itemize}
\end{frame}

% --- Slide 23: Cross-Embodiment Transfer ---
\begin{frame}{Cross-Embodiment Transfer}
\begin{center}
\includegraphics[width=0.75\textwidth,
  height=0.40\textheight,keepaspectratio]{%
  ../Figures/DreamZero/embodiment_transfer.pdf}
\end{center}

\begin{center}
{\small
\begin{tabular}{lc}
\toprule
\textbf{Method} & \textbf{Task Progress}\\
\midrule
DreamZero (baseline) & 38.3\%\\
+ Robot2Robot (YAM, 20 min)
  & \positive{55.4\%} (+44.6\% rel.)\\
+ Human2Robot (12 min)
  & \positive{54.3\%} (+41.8\% rel.)\\
\bottomrule
\end{tabular}
}
\end{center}
\end{frame}

% --- Slide 23b: Cross-Embodiment Commentary ---
\begin{frame}{Cross-Embodiment Transfer: Commentary}
\begin{softbox}
\higold{Commentary:} Video-only data (no action labels!)
opens a \hi{massive scaling pathway}.
Human video is orders of magnitude more abundant
than robot data.
\end{softbox}

\vspace{0.5em}
\begin{itemize}
    \item YouTube alone hosts \positive{$>$800M hours} of
          how-to and activity video
    \item The bottleneck shifts from ``collect robot data''
          to ``build a better IDM''
\end{itemize}
\end{frame}

% --- Slide 24: Few-shot ---
\begin{frame}{Few-shot Adaptation to New Robot}
\begin{center}
\includegraphics[width=0.82\textwidth,
  height=0.50\textheight,keepaspectratio]{%
  ../Figures/DreamZero/yam_transfer.pdf}
\end{center}

\begin{itemize}
    \item 55 trajectories ($\sim$30 min play data)
          on YAM robot
    \item Retains \positive{language following}
          + generalizes to novel objects
\end{itemize}
\end{frame}

% --- Slide 25: Free-form ---
\begin{frame}{Free-form Evaluation}
\begin{center}
\includegraphics[width=0.85\textwidth,
  height=0.55\textheight,keepaspectratio]{%
  ../Figures/DreamZero/free_form.pdf}
\end{center}

\hi{100+ novel tasks} via natural language instructions\\
``Pop the balloon'', ``Press elevator button'',
``Paint'', ``Iron clothes''
\end{frame}

% ==============================================================================
% PART 5: ABLATIONS & ANALYSIS
% ==============================================================================

% --- Slide 26: Ablation ---
\begin{frame}{Ablation: What Matters Most?}
\begin{center}
{\small
\begin{tabular}{llccc}
\toprule
& \textbf{Arch} & \textbf{Size} & \textbf{Data}
  & \textbf{Progress}\\
\midrule
\multicolumn{5}{l}{\textit{Q1.\ Data Diversity}} \\
& DZ (AR) & 14B & Repetitive
  & 33\% \\
& DZ (AR) & 14B & Diverse
  & \positive{50\%} \\
\midrule
\multicolumn{5}{l}{\textit{Q2.\ Model Scale}} \\
& DZ (AR) & 5B & Diverse & 21\% \\
& DZ (AR) & 14B & Diverse & \positive{50\%} \\
& VLA & 5B & Diverse & \negative{0\%} \\
& VLA & 14B & Diverse & \negative{0\%} \\
\midrule
\multicolumn{5}{l}{\textit{Q3.\ Architecture}} \\
& DZ (BD) & 14B & Diverse & 50\% \\
& DZ (AR) & 14B & Diverse & \positive{50\%} \\
\bottomrule
\end{tabular}
}
\end{center}
{\small All on PnP Easy, 50K steps, batch 32.
DZ = DreamZero.}
\end{frame}

% --- Slide 27: VLA vs WAM ---
\begin{frame}{Why VLAs Fail While WAMs Succeed}
\begin{columns}[T]
\column{0.48\textwidth}
\hi{VLAs:}
\begin{itemize}
    \item \negative{Static images}
          $\rightarrow$ no temporal dynamics
    \item Scaling does not help:
          \negative{0\% at 5B and 14B}
    \item Repetitive data + static priors
          = \negative{overfitting}
\end{itemize}

\column{0.48\textwidth}
\hi{WAMs:}
\begin{itemize}
    \item \positive{Video}
          $\rightarrow$ rich temporal dynamics
    \item Scaling improves video quality
          $\rightarrow$ \positive{better actions}
    \item Diverse data + video priors
          = \positive{generalization}
\end{itemize}
\end{columns}
\end{frame}

% --- Slide 27b: VLA vs WAM — Commentary ---
\begin{frame}{VLA vs.\ WAM: Commentary}
\begin{softbox}
\higold{Commentary:} A \hi{fundamental architectural
difference}, not data/compute. VLMs encode semantics;
video models encode physics. For robotics, physics wins.
\end{softbox}

\vspace{0.5em}
\begin{itemize}
    \item This echoes the ``bitter lesson'':
          the right \emph{representation}
          (video) matters more than scale
    \item Expect hybrid approaches that combine
          VLM semantics with WAM physics
\end{itemize}
\end{frame}

% --- Slide 28: Failure Analysis ---
\begin{frame}{Failure Analysis: When DreamZero Breaks}
\begin{center}
\includegraphics[width=0.75\textwidth,
  height=0.40\textheight,keepaspectratio]{%
  ../Figures/DreamZero/generated.pdf}
\end{center}
{\small Generated vs.\ executed trajectory pairs}

\vspace{0.3em}
\hi{Primary failure mode:}
\negative{video generation errors}
(hallucinations, wrong physics)
\end{frame}

% --- Slide 28b: Failure Causal Chain ---
\begin{frame}{Failure: Error Causal Chain}
\hi{Error causal chain:}
\begin{enumerate}
    \item Video backbone hallucinates
    \item Predicted frames show wrong trajectory
    \item IDM faithfully extracts actions from wrong video
    \item Robot executes incorrect motion
\end{enumerate}

\vspace{0.5em}
\begin{softbox}
\higold{Commentary:} The tight coupling is both strength
and weakness --- \hi{video quality is the bottleneck}.
\end{softbox}
\end{frame}

% ==============================================================================
% PART 6: CODE-LEVEL INSIGHTS
% ==============================================================================

% --- Section Divider ---
\begin{frame}[plain]
\begin{center}
\vfill
{\Huge\hi{Under the Hood}}\\[1em]
{\large\higold{Code-Level Insights}}
\vfill
\end{center}
\end{frame}

% --- Slide 29: Multi-Embodiment ---
\begin{frame}{Code: Multi-Embodiment Architecture}
\begin{methodbox}[Key Code Patterns]
\begin{itemize}
    \item \texttt{CategorySpecificLinear}:
          separate weights per embodiment
    \item \texttt{MultiEmbodimentActionEncoder}:
          category-specific MLPs
    \item Supports up to \hi{32 embodiments}
    \item Relative action normalization
\end{itemize}
\end{methodbox}

\vspace{0.5em}
\begin{softbox}
\higold{Commentary (from code):} Each robot gets its own
``translator'' while sharing the world model.
\end{softbox}
\end{frame}

% --- Slide 30a: Engineering ---
\begin{frame}{Code: Making 14B Real-Time}
\begin{methodbox}[Production-Grade Engineering]
\begin{itemize}
    \item \hi{VRAM management:} dynamic GPU--CPU offloading
    \item \hi{FP8 inference:} automatic scaling
    \item \hi{WebSocket server:} distributed multi-GPU
    \item \hi{DiT caching:} skip forward pass
          when velocity vectors converge
\end{itemize}
\end{methodbox}
\end{frame}

% --- Slide 30b: Quantization ---
\begin{frame}{Code: Quantization Stack}
\hi{Precision hierarchy:}
\begin{itemize}
    \item Weights + activations: NVFP4
    \item QKV, Softmax: FP8
    \item Non-linear ops: FP16
\end{itemize}

\vspace{0.3em}
\hi{Frame sequence:} 880 tokens
(33 frames $\times$ $\sim$27 tokens/frame)

\vspace{0.5em}
\begin{softbox}
\higold{Commentary (from code):} Far beyond typical
research code --- production-grade systems engineering.
\end{softbox}
\end{frame}

% --- Slide 31: Data Pipeline ---
\begin{frame}{Code: Data Pipeline}
\begin{methodbox}[Data Infrastructure]
\begin{itemize}
    \item \hi{LeRobot v2 format} with GEAR metadata
    \item Multi-view camera alignment
    \item Sharded dataset loading (10\% per shard)
    \item Action filtering: remove idle actions
\end{itemize}
\end{methodbox}

\vspace{0.5em}
\begin{softbox}
\higold{Commentary (from code):} Data infrastructure is
as important as the model architecture.
\end{softbox}
\end{frame}

% ==============================================================================
% PART 7: DISCUSSION & TAKEAWAYS
% ==============================================================================

% --- Section Divider ---
\begin{frame}[plain]
\begin{center}
\vfill
{\Huge\hi{What Does It All Mean?}}\\[1em]
{\large\higold{Discussion \& Takeaways}}
\vfill
\end{center}
\end{frame}

% --- Slide 32: Limitations ---
\begin{frame}{Limitations}
\hi{From the paper:}
\begin{enumerate}
    \item No established \negative{scaling laws} for WAMs
    \item Context window limited to \negative{6.6 seconds}
    \item Still \negative{7Hz} vs.\ 20Hz+ for VLAs
    \item Inherits behavior cloning limitations
          on \negative{sub-centimeter precision}
    \item Separate training per embodiment
\end{enumerate}
\end{frame}

% --- Slide 32b: Computational Requirements ---
\begin{frame}{Limitations: Computational Requirements}
\hi{Computational requirements:}
\begin{itemize}
    \item 2$\times$GB200 GPUs for real-time control
    \item Not yet feasible on consumer hardware
    \item 14B parameters = significant VRAM
\end{itemize}

\vspace{0.5em}
\hi{Long-horizon:}
\begin{itemize}
    \item 6.6s context $\rightarrow$ System 1 only
    \item Needs System 2 planner for extended tasks
\end{itemize}
\end{frame}

% --- Slide 33: Unsolved IDM ---
\begin{frame}{The Unsolved Inverse Dynamics Problem}
\hi{Adaptation is not free.}

\vspace{0.5em}
\hi{The credit card analogy:}
\begin{itemize}
    \item Dexterous hands: slide + pinch
    \item Grippers: push from edge
    \item \emph{Same task, different motions}
\end{itemize}

\vspace{0.5em}
\hi{IDM complexity scales with DOF:}
\begin{itemize}
    \item Simple grippers: relatively easy
    \item Dexterous hands: much harder
    \item Full humanoids: \negative{hardest}
\end{itemize}
\end{frame}

% --- Slide 34: Hardware vs Algorithm ---
\begin{frame}{The Hardware vs.\ Algorithm Race}
\begin{columns}[T]
\column{0.48\textwidth}
\hi{Hardware Path:}\\
If humanoid actuators mature quickly
$\rightarrow$ embodiment matching removes transfer loss

\vspace{0.5em}
\hi{Algorithm Path:}\\
If embodiment adaptation improves faster
$\rightarrow$ intermediate morphologies remain viable

\column{0.48\textwidth}
\begin{quotebox}
We are back to the age of research ---
fundamentals matter again.
\end{quotebox}
\hfill --- Jim Fan

\vspace{0.5em}
Both paths converge on the same data:\\
\higold{human video}
\end{columns}
\end{frame}

% --- Slide 35: Critical Assessment ---
\begin{frame}{Critical Assessment: Strengths}
\hi{Strengths:}
\begin{itemize}
    \item[\positive{\checkmark}] Comprehensive evaluation protocol
    \item[\positive{\checkmark}] Honest failure analysis
    \item[\positive{\checkmark}] Code and weights released
    \item[\positive{\checkmark}] Multiple embodiments tested
\end{itemize}

\vspace{0.5em}
\hi{Questions:}
\begin{itemize}
    \item How dependent on GB200 hardware?
    \item Will this work on consumer GPUs?
    \item Long-horizon tasks ($>$6.6s)?
\end{itemize}
\end{frame}

% --- Slide 35b: Critical Assessment — Missing ---
\begin{frame}{Critical Assessment: What's Missing}
\hi{Missing:}
\begin{itemize}
    \item Comparison with other recent WAMs
          (Cosmos Policy, UVAM) on same benchmarks
    \item ``From-pretrained'' baseline used
          \emph{different data}
    \item Multi-embodiment \emph{unified} training
\end{itemize}

\vspace{0.5em}
\begin{softbox}
\higold{Commentary:} Current zero-shot is ``AI Slop'' phase
(Joel Jang) --- correct direction, but
\hi{not deployment-ready} yet.
\end{softbox}
\end{frame}

% --- Slide 36a: Field Impact (1/2) ---
\begin{frame}{What DreamZero Means for the Field (1/2)}
\begin{keybox}[Personal Assessment]
\begin{enumerate}
    \item WAMs represent a \hi{paradigm shift}:
          from ``policy learning'' to ``world understanding''
    \item Jim Fan's ``second pre-training paradigm''
          is happening \higold{NOW}
    \item ``Diversity $>$ repetition'' could
          \positive{reshape robot data collection}
\end{enumerate}
\end{keybox}
\end{frame}

% --- Slide 36b: Field Impact (2/2) ---
\begin{frame}{What DreamZero Means for the Field (2/2)}
\begin{keybox}[Personal Assessment (cont.)]
\begin{enumerate}
    \setcounter{enumi}{3}
    \item Cross-embodiment transfer via video-only data
          is the \hi{sleeper result}
    \item Video backbone quality = policy quality\\
          $\rightarrow$ Watch for video model improvements
    \item ``Chain of thought in \hi{visual space}
          rather than language space''
\end{enumerate}
\end{keybox}
\end{frame}

% --- Slide 37: Key Takeaways ---
\begin{frame}{Key Takeaways}
\begin{enumerate}
    \item \hi{WAMs jointly predict video + action}
          $\rightarrow$ better generalization
    \item \hi{Diverse data + video priors}
          = zero-shot robot skills
    \item \hi{AR + teacher forcing + GT injection}
          = elegant inference
    \item \hi{38$\times$ speedup} makes 14B practical
    \item \hi{Cross-embodiment transfer}
          via video-only data
    \item \higold{``Train once on human experience,
          deploy everywhere''}
\end{enumerate}

\vspace{0.3em}
\begin{center}
\hi{The future of robotics is video-native,
not language-native.}
\end{center}
\end{frame}

% --- Slide 38: Thank You ---
\begin{frame}{Thank You \& References}
\begin{columns}[T]
\column{0.48\textwidth}
\hi{Paper \& Code:}
\begin{itemize}
    \item \citet{ye2026dreamzero}
    \item \texttt{dreamzero0.github.io}
\end{itemize}

\vspace{0.3em}
\hi{Blog Posts:}
\begin{itemize}
    \item Joel Jang: ``World Models for Robotics''
    \item Jim Fan: ``The Second Pre-training Paradigm''
\end{itemize}

\column{0.48\textwidth}
\hi{Key References:}
\begin{itemize}
    \item Wan2.1 \citep{wan2025wan}
    \item GR00T N1 \citep{bjorck2025gr00t}
    \item $\pi_0$/$\pi_{0.5}$
          \citep{black2024pi0,intelligence2025pi05}
    \item Flow matching \citep{liu2022flow,lipman2022flow}
    \item DROID \citep{khazatsky2024droid}
\end{itemize}
\end{columns}

\vspace{0.5em}
\begin{center}
\Large\hi{Q\&A}
\end{center}
\end{frame}

% --- Bibliography ---
\begin{frame}[allowframebreaks]{References}
\printbibliography[heading=none]
\end{frame}

\end{document}
